
El estudio de la paradoja EPR permitió comprender que el problema planteado por Einstein, Podolsky y Rosen no consistía en cuestionar directamente la validez de las predicciones de la Mecánica Cuántica, sino en discutir si la función de onda proporcionaba una descripción completa de la realidad física. Bajo su criterio de realidad y la suposición de localidad, EPR concluyó que dicha descripción era incompleta, dejando abierta la posibilidad de una teoría más detallada.

Sin embargo, mediante la deducción de la desigualdad original de Bell, pudimos observar que las suposiciones de realismo local imponen restricciones sobre las correlaciones entre los resultados de dos partículas. El enfoque gráfico utilizado permitió obtener estas restricciones sin necesidad de introducir explícitamente una distribución de variables ocultas. Al comparar la desigualdad obtenida con las predicciones de la Mecánica Cuántica para el estado singlete, encontramos que existen direcciones de medición para las cuales dicha desigualdad es violada. Por lo tanto, concluimos que no es posible reproducir todas estas correlaciones cuánticas mediante una teoría de variables ocultas locales que satisfaga las suposiciones consideradas.

Finalmente, este resultado también es relevante para el estudio del entrelazamiento atómico, pues el formalismo utilizado puede extenderse a sistemas atómicos efectivos de dos niveles preparados en estados entrelazados adecuados. Aunque nuestro trabajo se limitó al análisis teórico de las correlaciones del estado singlete, este desarrollo permite comprender cómo las desigualdades de Bell establecen restricciones sobre las posibles descripciones locales de los sistemas cuánticos.
